\section{Related Work}
\label{sec:related_work}

\paragraph{Self-correction without external feedback.}
Self-Refine~\cite{madaan2023selfrefine} uses one model in three roles (generator, critic and reviser) and iterates between them without any additional training. Reflexion~\cite{shinn2023reflexion} lets an agent write verbal reflections on the feedback it received for a task and keeps them in memory for later attempts. Both reported gains on a range of tasks, but later work questioned how much of the gain comes from the model itself. Huang et al.~\cite{huang2024cannot} showed that on GSM8K, CommonSenseQA and HotpotQA, part of the reported improvement came from using the ground-truth label to decide when to stop, and that purely intrinsic self-correction frequently made answers worse. Stechly et al.~\cite{stechly2023gpt4} studied iterative prompting of GPT-4 on graph colouring and found its self-critiques unreliable; most of the improvement came from a correct answer happening to appear among several candidates rather than from the critique itself. Kamoi et al.~\cite{kamoi2024when} reviewed the field and concluded that self-correction works reliably when trustworthy external feedback is available or when the model is fine-tuned for it, but that there is little evidence of success with prompting alone. Starting from this picture, we ask a narrower question: given a revision step that sometimes helps and sometimes hurts, can we decide in advance which answers to revise?

\paragraph{Confidence estimation for LLMs.}
A gate needs a confidence score, and several training-free ways to obtain one exist. The simplest is to ask the model directly. Xiong et al.~\cite{xiong2024can} evaluated such verbalised confidence together with sampling-based and hybrid methods across several models. They found that verbalised scores tend to be overconfident, and that checking consistency across multiple responses helps to reduce this. Self-consistency~\cite{wang2023selfconsistency} samples several reasoning paths and returns the majority answer; the share of samples that agree with a given answer is a simple measure of how sure the model is. Kadavath et al.~\cite{kadavath2022language} asked models to judge whether a proposed answer is true and read off the probability of ``True'', known as P(True); they found that larger models are well calibrated when questions are posed in a suitable format. We compare these three families of signals: verbalised, sampling-based and probability-based. Prior work mostly evaluates them as predictors of correctness, while we evaluate them as decision rules for whether to revise.

\paragraph{Confidence and self-correction.}
Two recent papers connect confidence directly to revision. Li et al.~\cite{li2024confidence} identify the model's confidence as a key hidden factor in intrinsic self-correction and propose the IoE prompt, in which the model assesses its own confidence and decides in the same prompt whether to keep or revise its answer. Yang et al.~\cite{yang2024confidence} define probabilistic measures for the confidence and critique abilities, report a trade-off between them under prompting, and propose a fine-tuning data format that improves both. The present study differs from these in three respects. First, the gate is an explicit threshold on a separately computed score rather than a decision left to the revision prompt, so different signals can be swapped in and compared under one protocol. Second, thresholds are selected on a development split and evaluated once on a test split. Third, an analytical model links the ROC curve of a signal to the accuracy gain it can deliver and explains when and why gating helps.
